\diarytitle{0043}{既知の1解 $y_1 = x$ を用いた階数低下法による第2解の導出}

2階線形同次方程式 $y'' + P(x)y' + Q(x)y = 0$ の1つの解 $y_{1}$（$y_1 \not\equiv 0$）が既知のとき、$y = y_{1}(x)\,v(x)$ とおいて代入すると、$y_1$ が解であることから $v$ を含まない項が消え、$v$ に関する方程式は $v''$ と $v'$ のみを含む形になる。すなわち $w = v'$ とおけば $w$ に関する1階線形方程式に帰着し、階数が1つ下がる。

$y_{1} = x$ が解であることは、$y_{1}' = 1$, $y_{1}'' = 0$ を代入して
$x^{2}\cdot 0 - 2x\cdot 1 + 2x = -2x + 2x = 0$
より確認できる。

$y = x\,v(x)$ とおくと
\[
y' = v + x v', \qquad y'' = 2v' + x v''
\]
これを方程式に代入すると
\[
x^{2}(2v' + x v'') - 2x(v + x v') + 2x v = 0
\]
\[
2x^{2} v' + x^{3} v'' - 2x v - 2x^{2} v' + 2x v = 0
\]
$-2x^2 v'$ と $2x^2 v'$ が打ち消し合い、$-2xv$ と $2xv$ も打ち消し合うので
\[
x^{3} v'' = 0
\quad\Longrightarrow\quad
v'' = 0 \quad (x \neq 0)
\]
よって $v' = C_{1}$、積分して $v = C_{1} x + C_{2}$。したがって
\[
y = x\,v = C_{1} x^{2} + C_{2} x
\]
$y_{1} = x$ に対応する項 $C_{2}x$ を除いた $y_{2} = x^{2}$ が第2の解であり、一般解は
\[
\boxed{\; y = C_{1} x^{2} + C_{2}\,x \;}
\qquad (C_{1},\ C_{2} \text{ は任意定数})
\]
（検算: $y_{2} = x^{2}$ について $y_{2}' = 2x$, $y_{2}'' = 2$ より $x^{2}\cdot 2 - 2x\cdot 2x + 2x^{2} = 2x^{2} - 4x^{2} + 2x^{2} = 0$。成立。$y_{1} = x$ と $y_{2} = x^{2}$ はロンスキアン $W = x\cdot 2x - 1\cdot x^{2} = x^{2} \neq 0\ (x\neq0)$ より1次独立。）
